\section{Inner edges and the two-band normal form}\label{sec:innerframes}

We now construct the data used in Definition~\ref{def:FPstar}.
The imported global reduction is the translated-well reduction of
\cite[Appendices A--D]{DMband}. We use its complex-energy transfer
bounds, physical Green kernels, Gram orthogonalization and localized
boundary weights. We do not import its Mathieu counting theorem as
a transport theorem. The changes needed at inner edges are proved
below. The lattice product estimates are recorded in
Appendix~\ref{app:libcalculus}.

\subsection{Two branches and the exact inverse Morse map}

Chambers' formula, in a common periodic Bloch gauge, is
\begin{equation}\label{eq:libchambers}
 \det(E-B_t(k))=P(E)-2\cos k-c(t),\qquad
 c(t)=2\Lambda\cos2\pi qt.
\end{equation}
This is the classical identity already used in
Section~\ref{sec:sharp}; see \cite{Chambers1965} and
\cite[Methods and conventions]{DMband}.
The phase-union bands are the intervals on which $|P|\le R_*$.
An inner gap $\gamma$, possibly closed, has edges $E_a\le E_a'$ and
one critical point $c_\gamma$. Write
\[
 P(E_a)=P(E_a')=sR_*,\quad
 P(c_\gamma)=s(R_*+\mathrm{ex}),\quad s\in\{-1,1\},\quad
 \mathrm{ex}\ge0.
\]
The turning phase satisfies $c(t_a)=2s\Lambda$ and
$\cos k_a=s$. Roots of $P+sR_*$ are counted with multiplicity.
Their real-rootedness follows by realizing the polynomial as a
Hermitian fibre characteristic polynomial.

\begin{lemma}[Curvature and the Morse chart]\label{lem:libMorse}
Let $d_S=4\lambda+5$ and $M_j=\sup|P^{(j)}|$ in the fixed complex
neighborhood used for the rational derivative bounds. For large $q$,
\begin{align}
 |P''(c_\gamma)|&=(2R_*+\mathrm{ex})
       \sum_{f:P(f)=-sR_*}(c_\gamma-f)^{-2}
       \ge 2R_*q/d_S^2,\label{eq:libcurvature}\\
 |P'(E_a)|&\le M_2|E_a-c_\gamma|,\qquad
 \mathrm{ex}\ge |P'(E_a)|^2/(2M_2).\label{eq:libexcess}
\end{align}
There is $\rho_c=e^{-d_hq}/(C_0q^2)$, with $C_0$ sufficiently large,
on a contracted disc about $c_\gamma$ on which
\begin{equation}\label{eq:libMorse}
 P(E)=s(R_*+\mathrm{ex})-s\zeta(E)^2,\quad
 \zeta(c_\gamma)=0,\quad z_\gamma:=\zeta'(c_\gamma)
      =\sqrt{|P''(c_\gamma)|/2}>0,
\end{equation}
whenever the two edges lie in its inner quarter.
The map $\zeta$ is analytic and univalent there, is real on the real
axis, and $\zeta'(E)/z_\gamma=1+O(1/q)$.
\end{lemma}

\begin{proof}
Put $Q=P+sR_*$. On the root-free interval containing the gap,
$(Q'/Q)'=-\sum_f(E-f)^{-2}$. At $c_\gamma$ the first logarithmic
derivative vanishes and $Q(c_\gamma)=s(2R_*+\mathrm{ex})$.
This proves~\eqref{eq:libcurvature}. Integrating the same identity
also gives $|P'(E_a)|\ge2R_*q|c_\gamma-E_a|/d_S^2$.
The mean-value theorem gives the first inequality
in~\eqref{eq:libexcess}. Integrate
$|P'(E)|\ge |P'(E_a)|-M_2|E-E_a|$ up to its zero, which lies before
$c_\gamma$, to get the second.

The derivative bounds below are used only in the fixed complex neighborhood of the rational phase-union bands. To justify their use at a narrow-gap center, integrate the logarithmic derivative once more:
\[
 \log\left(1+\frac{\mathrm{ex}}{2R_*}\right)
 \ge\frac{q|E_a-c_\gamma|^2}{2d_S^2}.
\]
Thus the center of an N gap tends to its edge and belongs to that neighborhood. The triangular integral proving~\eqref{eq:libexcess} needs only an $O(1/q)$ segment at the edge. If the critical point is farther away, take the same segment and enlarge the fixed derivative constants. No growth bound at a far gap center is assumed.

The rational derivative bounds are
$M_j\le C_{j,h}q^j e^{hq}$, by
\cite[Proposition 3.3(2)]{DMamo} and
\cite[Appendix A, Polynomial derivatives and the top contour]{DMband}.
Taylor-expand at $c_\gamma$ and take the analytic square root of the
nonzero quadratic factor. Bound that factor and its derivative on
a larger disc before contracting to $\rho_c$. Increasing $C_0$
makes its derivative $O(1/q)$ relative to its value at the center.
This proves univalence and all assertions about $\zeta$.
\end{proof}

On an isolated two-band chart let $h(k,t)$ be the frozen restriction
of $B_t(k)$ to these two bands, in an orthonormal analytic frame.
Put
\[
 T_s(k)=2-2s\cos k,\quad U(t)=2\Lambda-sc(t),\quad
 x=\mathrm{ex}+T_s(k)+U(t).
\]
Cayley--Hamilton and~\eqref{eq:libchambers} give the exact identity
\begin{equation}\label{eq:libDirac}
 Z=\zeta(h),\qquad Z^2=xI_2,\qquad \tr Z=0.
\end{equation}
Indeed the two distinct energies in the chart have opposite
$\zeta$-values. Continuation proves the trace identity at a touching
point as well. For real $k,t$, $Z$ is Hermitian.

Write $c=c_\gamma$ and keep the full inverse
\[
 \psi(w)=z_\gamma\bigl(\zeta^{-1}(w)-c\bigr),\qquad Y=\psi(Z).
\]
Define analytic scalar functions
\begin{equation}\label{eq:libinverse}
 \vartheta(x)=\frac{\psi(\sqrt x)+\psi(-\sqrt x)}2,\qquad
 \Omega(x)=\left(\frac{\psi(\sqrt x)-\psi(-\sqrt x)}2\right)^2.
\end{equation}
They are analytic also at $x=0$. With
$\widetilde Y=Y-\vartheta(x)I_2$, one has
\begin{equation}\label{eq:libfactor}
 \widetilde Y^2=\Omega(x)I_2,\qquad
 (Y-Y_0)(Y+Y_0-2\vartheta(x))=\Phi(x)I_2,
 \quad \Phi(x)=\Omega(x)-(Y_0-\vartheta(x))^2.
\end{equation}
For admissible $Y_0$ in the contracted chart,
$\Phi'=1+O(1/q)>0$. The inverse estimates imply this after contracting
the chart and choosing $\xi$ larger than its fixed losses.
These are exact frozen identities. In general $Y\ne Z$.
Replacing the inverse by its tangent would already introduce a
nonzero error at $\delta=0$.

\subsection{The wide and narrow classes}

Use the final threshold
\begin{equation}\label{eq:libclasses}
 \mathrm{ex}_*=\Lambda e^{-4d_hq}/q.
 \qquad \mathrm W:\ \mathrm{ex}>\mathrm{ex}_*,\qquad
 \mathrm N:\ \mathrm{ex}\le\mathrm{ex}_*.
\end{equation}
Class W uses separate one-band frames at its two edges. Class N
uses one frame containing both bands. Every exactly closed pair
belongs to N. The threshold is an eventual asymptotic threshold,
not a classification inferred from a finite numerical sample.

\begin{lemma}[Isolation and localized frames]\label{lem:libframes}
For every sufficiently small $\xi>0$, choose $d_h>0$ sufficiently
small after all fixed derivative orders. For large $q$ the following
hold uniformly on $|t-t_a|\le t_{\rm fr}$.
In W each edge has a one-root contour of radius at least
$e^{-6d_hq}/q^C$, with discriminant on its boundary at least
$\Lambda e^{-11d_hq}/q^C$.
In N the Morse contour contains exactly two roots, with boundary
discriminant at least $\Lambda e^{-2d_hq}/q^C$.
Both give Wannier frames with, for the finite derivative orders used
here,
\begin{equation}\label{eq:libtails}
 |\partial_t^j w_{t,\ell}(x)|\le
 \mathcal R_q e^{-(L-4\tau)\dist(x,\mathcal A)},
 \qquad \#\mathcal A\le2.
\end{equation}
Here $\mathcal A$ is a finite set of lifted anchors in one cell.
The translated global synthesis and its orthogonalization obey
physical weighted row and column Schur bounds, retaining an
interwell error at most $\mathcal R_q e^{-c/|\delta|}$.
\end{lemma}

\begin{proof}
In W, if $c_\gamma$ is outside a fixed edge neighborhood, the logarithmic-derivative lower bound gives $|P'(E_a)|\ge c_\lambda q\Lambda$. Otherwise the segment from $E_a$ to $c_\gamma$ is inside the derivative neighborhood, and $\mathrm{ex}\le M_1|c_\gamma-E_a|$ combines with that lower bound. In either case the threshold~\eqref{eq:libclasses} gives
$|P'(E_a)|\ge\Lambda e^{-5d_hq}/q^C$.
Take a fixed contraction of
$\min\{|P'(E_a)|/(2M_2),R_*/M_1\}$ as the one-root radius.
Taylor's theorem and Rouch\'e give the stated contour estimates.
In N,~\eqref{eq:libcurvature} and~\eqref{eq:libclasses} put both roots
in the inner quarter of the Morse disc for large $q$.
Rouch\'e counts exactly two roots there. The phase drift on the
frame chart is $O(\Lambda e^{-\xi q/2})$ and is smaller than the
boundary discriminant. Contract the complex Bloch strip accordingly.

The periodic Green-function argument in
\cite[Appendix B, Periodic Green function and Analytic fibre data]{DMband}
uses only this large discriminant and the complex transfer bounds.
It therefore applies on these contours. It gives cyclic Green
kernels and analytic fibre data on a strip of width
$(L-O(d_h)-\text{chosen loss})q$.
Choose a nondegenerate anchor Gram matrix at the reference real
fibre. In rank two, pivoted Cholesky gives determinant at least
$2/q^2$; the overlap anchors in Lemma~\ref{lem:libanchors} will give
another valid choice at an open N gap.
For its complex continuation use the periodic trace and determinant
of the Gram matrix. Endpoint phases cancel in principal minors.
Their Fourier coefficients are bounded on the large strip, so their
variation on a contracted strip is small. The identity
\[
 G^{1/2}=\frac{G+\sqrt{\det G}\,I_2}
                 {\sqrt{\tr G+2\sqrt{\det G}}}
\]
continues the inverse square root holomorphically from the positive
real reference matrix. This does not assert entrywise Fourier
flatness in an incompatible gauge.

Include the anchor endpoint phases before inverse Bloch transform,
and choose a common periodic gauge for both columns. Place their
finite lifts along the shorter arc, of length $r\le q/2$.
Contour shifting gives~\eqref{eq:libtails}. The tail of one Wannier
function is measured from this finite set, rather than from its
periodic repetition. Translating each well to the reference chart
and applying the physical Schur orthogonalization of
\cite[Appendix C, Overlap cancellation and global orthogonalization]{DMband}
gives the global assertions. Across wells sum over distinct physical
centers; the number of anchors per center is bounded. This proves
both Schur sums without summing a nonzero error once per well.
\end{proof}

The raw global construction remains valid when $b\ge2\xi$.
Only its freezing comparison uses a lower bound on $b$.
On patches of $\ell_c=\lceil(q|\delta|)^{-1/2}\rceil$ cells the phase
variation is $O(\sqrt{|\delta|/q})$. It is smaller than the frame
chart and than the complementary energy margin below.
The estimates normalized by a level spacing need not be small on
these scales. We keep the raw estimates and enlarge $D$ instead.

\begin{proposition}[Compression and direct complement exclusion]
\label{prop:libcompression}
For $2\xi\le b\le L-\xi$, the global frames above have a self-adjoint
effective operator $h_s$ and, throughout their contracted class
windows,
\begin{equation}\label{eq:libcompdata}
 \|(QHQ-E)v\|\ge m\|v\|\quad(v\in Q\ell^2),\qquad
 m\ge q^{-C}e^{-(\xi/2+Cd_h)q}.
\end{equation}
The bulk coupling is $c_{\rm bulk}\le|\delta|\mathcal R_q$.
In N the safe compression error and separate end coupling are
\begin{equation}\label{eq:libranktwoerrors}
 \gamma_{\rm cmp}\le |\delta|\mathcal R_q
       e^{-(L-4\tau)(q-r)}+\delta^2\mathcal R_q
       +\mathcal R_q e^{-c/|\delta|},\qquad
 c_{\rm end}\le e^{-(L/2-\text{chosen loss})q}.
\end{equation}
The end leg is localized in the end strips, with decaying Fourier
and interwell tails. In W the corresponding one-anchor exponent
is retained, together with the positive boundary weight of
\cite[Appendix D, Interior coupling and omitted translates]{DMband}.
\end{proposition}

\begin{proof}
The translated matrix expansion, its Hermitian compression, and the
separation of bulk and end couplings are the kernel calculations in
\cite[Appendices C--D]{DMband}. They use nearest-neighbor hopping,
the physical localization kernels and finitely many phase
derivatives. With two anchors the convolution distance for a
nonzero cell displacement is $|m|q-r$, rather than $|m|q$.
This gives~\eqref{eq:libranktwoerrors}, with both factors at an end
interface retained. It also proves the differentiated seminorms
in Appendix~\ref{app:libcalculus}. A column-frozen matrix is replaced
by its Hermitian part before any comparison.

For complement exclusion the omitted frozen branches are at least
\[
 g_{\mathrm W}\ge\min\{\rho_a/2,M''/(C|P'(E_a)|)\},\quad
 g_{\mathrm N}\ge\min\{\rho_c/2,
       M''/[C\sqrt{M_2(\mathrm{ex}+M'')}]\},
 \quad M''=c_\lambda\Lambda e^{-\xi q/2}
\]
from the class window. Inside a frame patch this follows from the
one-root or two-root contour. Outside it, the phase value is
separated from the turning value by $M''$; the inverse derivative
or Morse divided difference gives these bounds. Use a common lower
bound $g\ge q^{-C}e^{-(\xi/2+Cd_h)q}$ in all patches.

Put $\eta_{\rm fr}=C\sqrt{|\delta|/q}$. On each patch let
$A_a=B_a-E$ be the frozen operator and $\Pi_a$ its retained one- or
two-band projection. The common margin $g$ gives
$\|A_a f\|\ge g\|(1-\Pi_a)f\|$. Freezing and
$(x-y)_+^2\ge x^2/2-y^2$ imply
\[
 \|(H-E)f\|^2\ge\frac{g^2}{2}
       \bigl(\|f\|^2-\|\Pi_a f\|^2\bigr)
       -\eta_{\rm fr}^2\|f\|^2.
\]
For $v\in Q\ell^2$, apply the genuine IMS formula to $(H-E)^2$
and sum with $f=\chi_a v$. Its error is $C|\delta|/q$.
The supported joint Bessel estimate of
\cite[Appendix D, Excluding the complementary subspace]{DMband},
summed over the at most two anchors, is
\[
 \sum_a\|\Pi_a\chi_a v\|^2\le
       \mathcal R_q|\delta|q^{-1}\|v\|^2.
\]
The patch-support factors are included in both Schur sums. Hence
\[
 \|(H-E)v\|^2\ge
 \left[\frac{g^2}{2}-\frac{g^2}{2}\mathcal R_q\frac{|\delta|}{q}
                  -C\frac{|\delta|}{q}\right]\|v\|^2.
\]
Since $b\ge2\xi$ and $d_h$ is chosen after the fixed constants,
$\mathcal R_q|\delta|/q=o(1)$ and $|\delta|/q=o(g^2)$.
Therefore $\|(H-E)v\|\ge(g/\sqrt2-o(g))\|v\|$.
The complete coupling, including the conservative rank-two end
term, is $o(g)$. Consequently
\[
 \|(QHQ-E)v\|\ge\|(H-E)v\|-\|\Pi HQ\|\|v\|
       \ge(g/2)\|v\|.
\]
This proves the gap for the compression directly. A spectral gap of
$H$ alone would not imply it. Both nearby branches have been
included in N, also when their gap is positive but very small.
\end{proof}

\subsection{The curvature dichotomy and overlap anchors}

\begin{lemma}[Overlap anchors at every open N gap]\label{lem:libanchors}
At an open class-N gap, let $u,u'$ be the normalized real edge
vectors at $(t_a,k_a)$ and let $\mathcal R$ be the signed reflection
commuting with that fibre. There is a non-fixed site $x_*$ such that
\begin{equation}\label{eq:libanchor}
 |u(x_*)u'(x_*)|\ge q^{-C}e^{-Cd_hq},\qquad
 \det G_{\{x_*,\mathcal Rx_*\}}=4u(x_*)^2u'(x_*)^2.
\end{equation}
If $r=\dist_{\Z/q\Z}(x_*,\mathcal Rx_*)\le q/2$, then
\begin{equation}\label{eq:libtunnel}
 |P'(E_a)|\le e^{hq}\mathcal R_q e^{-(L-2\tau)r}.
\end{equation}
\end{lemma}

\begin{proof}
Let $V'=\partial_t B_{t_a}$ and
$M=\langle u',V'u\rangle$. Twice differentiating Chambers' identity,
then using simple-eigenvalue perturbation theory, gives
\begin{equation}\label{eq:libRellich}
 \frac{c''(t_a)}{P'(E_a)}=
 \langle u,V''u\rangle+
 \frac{2|M|^2}{E_a-E_a'}+
 2\sum_{j\ne a,a'}\frac{|\langle u_j,V'u\rangle|^2}{E_a-E_j},
 \qquad |c''(t_a)|=8\pi^2q^2\Lambda.
\end{equation}
The derivative in this identity is signed. Other degenerate
subspaces are summed by their projectors. A whole phase-union band
separates the pair from every other same-sign root, so the other
energy denominators have modulus at least $2R_*/M_1$.
Also $\|V'\|\le4\pi\lambda$ and $\|V''\|\le8\pi^2\lambda$.

If the partner term has modulus less than $|E_a''|/3$, the other
terms imply $|P'(E_a)|\ge c_\lambda q\Lambda e^{-d_hq}$.
Equation~\eqref{eq:libexcess} then gives
$\mathrm{ex}\ge c_\lambda'\Lambda e^{-3d_hq}$.
For large $q$ this is W under~\eqref{eq:libclasses}.
Otherwise
\[
 |M|^2\ge \frac{4\pi^2q^2\Lambda}{3M_2}
       \ge c_\lambda e^{-d_hq}.
\]
Every open N gap is therefore in this second alternative.
Reflection anticommutes with $V'$, so its two edge vectors have
opposite reflection parities. Since
$|M|\le4\pi\lambda q\max_x|u(x)u'(x)|$, a site $x_*$ satisfies
\eqref{eq:libanchor}. At a reflection-fixed site one of these
opposite-parity vectors vanishes. Hence $x_*$ is not fixed.
The Gram determinant follows by writing the two reflected columns
in the edge basis.

Let $\Pi_\gamma$ be the frozen pair projector and $ed\in\{-1,1\}$
the reflection sign at $x_*$. The exact two-dimensional pencil
identity is
\[
 (E_a-E_a')(\Pi_\gamma)_{x_*x_*}
 =2ed(B\Pi_\gamma)_{x_*,\mathcal Rx_*}
   -ed(E_a+E_a')(\Pi_\gamma)_{x_*,\mathcal Rx_*}.
\]
The two off-diagonal kernels decay in the cyclic distance $r$ by
the contour Green bounds. The diagonal is at least
$2|u(x_*)u'(x_*)|$. Thus
$|E_a-E_a'|\le\mathcal R_q e^{-(L-2\tau)r}$.
Finally $|P'(E_a)|\le M_2|E_a-E_a'|$, which proves
\eqref{eq:libtunnel}. No maximizer of a projector diagonal has been
assumed non-fixed.
\end{proof}

At an exactly closed pair we use the rank-two Riesz projector;
\eqref{eq:libRellich} is not asserted for an individual vector at a
double eigenvalue. For closed and near-closed pairs the general
pivot anchors in Lemma~\ref{lem:libframes} give the longer-arc tail
with $r\le q/2$. In the shallow comparison it suffices to use the
cutoff $\sqrt{\mathrm{ex}}\le e^{(\xi/2+6\tau)q}$.
Open N gaps outside that cutoff use Lemma~\ref{lem:libanchors}.
At the longer-arc jump bond both endpoints are included. If
\begin{equation}\label{eq:libF}
 F=\mathcal R_q e^{-(L-4\tau)(q-r)},
\end{equation}
then $\|\un_{\rm bonds}\Pi\|^2\le F$.
This is a physical operator bound from row and column Schur sums,
rather than a single frozen projector diagonal.

\subsection{The one-band W package}

For a W edge put $P_a'=|P'(E_a)|$ and
$\sigma_{\rm out}=s\,\operatorname{sgn}P'(E_a)$.
The positive inward depth operator is
$P_a'\sigma_{\rm out}(E_a+s_a(t_a)-h_s)$.
Here $s_a$ is the orbital term, including the small parameter.
The inverse branch $E(y)$ satisfies $P(E(y))=sR_*-sy$, and
\[
 F_a(y)=P_a'|E(y)-E_a|,\qquad
 |F_a'(y)-1|\le4M_2y/(P_a')^2
       \le q^C e^{-(\xi-Cd_h)q}
\]
on $0\le y\le C\Lambda e^{-\xi q}$.
Taylor's theorem proves this after contracting the inverse domain.
The orbital correction is the classical Rammal--Wilkinson term
\cite{Wilkinson1987,RammalBellissard,HSII}; its uniform local compression is
\cite[Appendix C]{DMband}.
A possible odd part has size
$\mathcal R_q\sqrt\kappa\sqrt U$, without an extra $|\delta|$.
Its first two phase derivatives move the minimum by
$\mathcal R_q$ cells and preserve quadratic growth there.
Thus a margin $D$ gives one shell of width
$C\sqrt{D/\kappa}+\mathcal R_q$ near the minimum, or two shells of
width $CD/\sqrt{\kappa e}+\mathcal R_q$ away from it.
The coarse unit is
\[
 A_a=D_a/P_a',\qquad
 D_a=C\{1+P_a'(\gamma_{\rm cmp}+c_{\rm bulk}^2/m
                                      +\|W_{\rm bd}\|)\}.
\]
Diagonal dominance after the shell penalty gives a gap $3A_a$.
The one-anchor midpoint tail, the weight estimates of
Appendix~\ref{app:liblocal}, and Proposition~\ref{prop:libcompression}
prove the W package. At low $b$ the raw estimate gives
\[
 \sqrt{D_a/\kappa}\le
 \mathcal R_q\bigl(\kappa^{-1/2}+e^{\xi q/4}\bigr).
\]
The padding $e^{(\xi/2+12\tau)q}$ absorbs the additional term.
Top and bottom use the same one-band argument and the same coarse
unit. These estimates are uniform at all depths of their class.
The remaining N shell and weighted estimates are proved next.
